%!TEX program = uplatex
\RequirePackage{plautopatch}
\documentclass[uplatex,dvipdfmx]{jlreq}
\usepackage{amsmath,amssymb}

\pagestyle{empty}

\newcommand{\Diff}{\mathrm{Diff}^r_+(S^1)}
\newcommand{\R}{\mathbf{R}}

\begin{document}

%======================================================================
% 表紙
%======================================================================
\begin{center}
  \vspace*{20pt}
  {\Huge\bfseries 実１次元力学系}

  \vspace{30pt}

  {\Large\bfseries 群の円周への作用をそのダイナミックスから解析する。}

  \vspace{50pt}

  {\large 前文}

  \vspace{10pt}

  \parbox{0.85\textwidth}{
    \large\bfseries
    常微分方程式の定性的理論から発生した円周の同相写像の
    力学系的研究は様々な群の円周への作用を研究する
    大きな研究分野へと育っている。多くの興味深い群作用を例示し、
    群作用の共役不変量から群作用の分類に迫る。
  }
\end{center}

\vfill


\newpage

%======================================================================
% 講演１：坪井俊
%======================================================================
\begin{center}
  {\large\bfseries 実１次元力学系}
\end{center}
\begin{flushright}
  \textbf{坪井　俊（東大数理）}
\end{flushright}

\smallskip

本講演では実１次元力学系入門としていくつかのトピックの解説をおこなう。
３回の60分講演で、次のような内容を予定している。例を豊富に与えることとし、証明は与えない。

\bigskip
\noindent
{\large\textbf{第１回}}

\medskip
\noindent
\textbf{１次元力学系研究の由来。円周の同相写像を研究することとは何か。}

\medskip

常微分方程式の解の定性的研究はポアンカレに始まるといわれている。
時刻によらない常微分方程式がたくさんの第一積分を持てば、
解はある程度低い次元の空間にとどまる。
この空間がコンパクトで１次元ならば解は円周上を動き、理解しやすい対象である。
この空間がコンパクト２次元で、常微分方程式が定数解を持たないとすると、
２次元多様体はトーラス（またはクラインボトル）であり、
解曲線に横断的な円周が存在する。
さらに、円周と２度以上交わる解曲線が存在すれば、
すべての解曲線は円周と交わり、このような解曲線の
回帰写像として円周の同相写像が定まる。
この同相写像の軌道の様子を調べれば常微分方程式の解曲線の様子がわかる。
円周の同相写像のイテレーションに周期点があれば周期解があり、
イテレーションの軌道が稠密なら、解曲線はトーラス上で稠密となる。

\medskip
\noindent
\textbf{円周の同相写像の表示、同相写像のイテレーションの様子。
回転数とは何か。ダンジョワの定理。}

\medskip

円周の向きを保つ同相写像は、直線の「周期的な」同相写像に持ち上げられる。
この周期的な同相写像のグラフを書くこと、そのイテレーションを書くことは
現在ではパソコン上で容易にできる。
イテレーションの様子を見ると豊富なダイナミックスが含まれていることがわかる。
直線の周期的な同相写像の共役不変量が「移動数」であり、
これが円周の同相写像の共役不変量である回転数 $\in \mathbf{R}/\mathbf{Z}$ を与える。
ダンジョワによれば $C^2$ の微分同相写像は、回転数が無理数ならば、
その回転数の回転と位相共役である。
微分共役の条件は難しい。アーノルド、エルマン、ヨコス達の結果がある。
また、$C^2$ よりも微分可能性が低い場合もダンジョワ、ピクストン、$\cdots$ の仕事がある。
（回転数が有理数ならば周期軌道が存在する。）
この回転数は $\mathrm{Homeo}(S^1)$ の２次元有界コホモロジーにある有界オイラー類の
$\mathbf{Z}$ の２次元有界コホモロジー（$\cong \mathbf{R}/\mathbf{Z}$）への引き戻しと理解される。
（98年のＩＣＭでのマクマレンの講演の中に微分可能な同相写像で、
臨界点を持つものが扱われていたことを付記しておく。）

\medskip
\noindent
\textbf{円周の同相群の有限生成部分群、軌道、極小集合。サックステーダーの定理。}

\medskip

軌道の概念は円周の同相群の有限生成部分群に対し定義される。
この有限生成部分群の作用に対し、空でない不変な閉集合のうち極小のものを極小集合とよぶ。
極小集合は有限個の点、円周全体、またはカントール集合となる
（最後の場合を例外極小集合と呼ぶ）。
どのような極小集合を持っているかは有限生成部分群の作用のダイナミックスにとって重要な不変量である。
サックステーダーによれば、有限生成部分群 $G$ の $C^2$ 級の作用が例外極小集合を持てば、
例外極小集合に含まれる可算個の点 $p$ で、$p$ を固定点とする $g \in G$ が存在し、
$g$ の $p$ での微分は１と異なる。すなわち $p$ は線形にアトラクティブまたはリパルシブとなる。

\bigskip
\noindent
{\large\textbf{第２回}}

\medskip
\noindent
\textbf{円周の同相群の有限生成部分群の例。その力学系的性質。
可換部分群。可解部分群。
有限生成群の作用の軌道の増大度。多項式増大度を持つ作用。}

\medskip

$SO(2)$ の無理数回転と可換な同相写像は $SO(2)$ の回転に限る。
皆川によれば、$S^1$ の区分線形同相写像の群 $PL(S^1)$ に
$SO(2)$ と位相共役なものが存在する。
これは皆川氏の講演でも述べられると思うが驚くべき現象である。

周期軌道をもつ同相写像は何回かイテレートして固定点を持つようにする。
固定点をもつ同相写像 $f$ と可換な同相写像 $g$ については、
固定点が共通である場合が最も面白い。
そのような場合、コペル、タケンズの結果を使うと、
$f$, $g$ が実解析的で階数２の自由加群を生成するならば、
$f$, $g$ ともに円周上のベクトル場 $\xi$ の積分となることがわかる。

$G = \langle f, g \mid fgf^{-1}g^{-1} = g^{-1}fgf^{-1} \rangle$
という群の作用は、とくに $f$, $g$ の作用が可換とならず、
ともに固定点を持つ場合に色々と興味深い例を与える。
実解析的な作用は可解群を経由し、基本的に上半三角行列
$\subset PSL(2, \mathbf{R}) \subset \mathrm{Diffeo}(S^1)$
に含まれるものとなる。

$G_1 = \langle f, g \mid f^k g f^{-k} \text{ と } g \text{ は可換},\ k \in \mathbf{Z} \rangle$
を考えると、$G_1$ の $C^\infty$ 級の作用のなかに、$g$ の固定点集合が $f$ の固定点集合を
真に含むことがある。このなかにステアケースと呼ばれる作用がある。
それは $f$ の固定点を両端とする区間の中の $f$ の基本領域に $g$ の台があるような場合である。
$g$ の台の内点 $x$ の軌道は $\mathbf{Z}^2$ を辞書式にならべた形であり、
$G$ の長さ $n$ 以下のワードによる $x$ の像の個数は $n^2$ のオーダーで増える。
このようなものを一般化して多項式増大度をもつ軌道を持つ有限生成群の作用も構成される。
すべての軌道が多項式増大度をもつような作用は大まかな分類が可能である。

群自身の増大度の問題はグロモフ、グリゴルチュク $\cdots$ により研究されているが、結構難しい。

区分線形同相写像の群 $PL(S^1)$ については、
皆川氏の講演でその豊かさが実感されていると思うが、
$PL(S^1)$ の部分群にはヒグマン・トンプソンの群と呼ばれる有限表示無限単純群がある。
たとえば、
\[
  \langle x_0, x_1 \mid
  x_1^{-1} x_0^{-1} x_1 x_0 x_1 = x_0^{-2} x_1 x_0^2,\;
  x_0^{-3} x_1 x_0^3 = x_1^{-1} x_0^{-2} x_1 x_0^2 x_1 \rangle
\]
がそのような群である。実際、$PL(S^1)$ は実例の宝庫である。

\bigskip
\noindent
{\large\textbf{第３回}}

\medskip
\noindent
\textbf{指数関数的増大度を持つ群。フックス群。
群作用のエントロピー。ゴドビヨン・ヴェイ不変量。
ジス・ランジュヴァン・ワルチャックの定理とドゥミニーの定理。}

\medskip

一般の有限生成群の作用を考えると指数関数的増大度をもつ軌道の存在が予想される。
このような作用の最も重要な例は何度も登場しているフックス群の作用である。
有限表示群に対しては、ボーエン、ジス・ランジュヴァン・ワルチャックにより定義された
群作用のエントロピーという量がある。
$S^1$ の部分集合 $A$ について、その任意の２点が長さ $n$ 以下のワードの作用で
$\varepsilon$ 以上離れることの出来るとき、$A$ は $(n,\varepsilon)$ 分離的ということにする。
$N(n,\varepsilon)$ で $(n,\varepsilon)$ 分離的部分集合の元の最大個数をあらわす。
$N(n,\varepsilon)$ はもちろん $\frac{1}{n}$ 以上であるが、横断的増大度を表すもので、
作用が複雑なら指数関数的にも増えうる。（この量についての江頭による研究がある。）
\[
  \lim_{\varepsilon \to +0} \limsup_{n \to \infty} \frac{1}{n} \log N(n,\varepsilon)
\]
をエントロピーと呼ぶ。
ジス・ランジュヴァン・ワルチャックによれば、
群作用のエントロピーが正であることと弾性軌道の存在は同値である。
$x$ の軌道が弾性軌道であるとは $x$ の軌道上の点 $g_1(x)$ が $x$ を
アトラクティブな固定点とする $g_2$ のアトラクティングな領域に含まれることをいい、
この場合いわゆるピンポンの議論により $x$ の軌道は指数関数的増大度を持つ。
曲面群の円周への $C^2$ 級作用に対しては、ゴドビヨン・ヴェイ不変量が
曲面群の２次元コホモロジーの元として定義される。
ドゥミニーによればゴドビヨン・ヴェイ不変量が零でなければ弾性軌道が存在する。
従って、曲面群の円周への作用のゴドビヨン・ヴェイ不変量が零でなければ、
群作用のエントロピーは正である。

\medskip
\noindent
\textbf{収束群についてのツキア・ガバイ・キャッソンの定理。}

\medskip

$\mathrm{Homeo}(S^1)$ の部分群 $G$ が収束群とは、
$G$ の任意の無限列 $g_i$ に対し、部分列 $f_i$ と点 $x$, $y$ を選んで、
$f_i$, $(f_i)^{-1}$ が $S^1 \setminus \{x, y\}$ 上広義一様に定値写像 $x$, $y$ に収束する
ように出来ることをいう。
ツキア・ガバイ・キャッソンにより、このような群は $PSL(2,\mathbf{R})$ の離散部分群となることが示されている。
これはフックス群の力学系理論的な特徴付けとして重要なものである。
余裕があればこの話題にも触れたい。

\medskip

\begin{center}
  {\large 参考文献}
\end{center}

\begin{itemize}
  \setlength{\itemsep}{3pt}
  \item 青木統夫，力学系・カオス，共立出版（1996）．
  \item A.~Casson and D.~Jungreis,
        % Convergence groups and Seifert fibered 3-manifolds.
        \textit{Invent.\ Math.}\ \textbf{118} (1994), no.~3, 441--456.
  \item コディントン・レヴィンソン，常微分方程式論（下），吉田節三訳，吉岡書店（1969）．
  \item D.~Gabai,
        % Convergence groups are Fuchsian groups.
        \textit{Ann.\ of Math.}\ (2) \textbf{136} (1992), no.~3, 447--510.
  \item E.~Ghys,
        % Groupes d'homeomorphismes du cercle et cohomologie bornee.
        \textit{Contemp.\ Math.}\ \textbf{58}, III, Amer.\ Math.\ Soc.\ (1987), 81--106.
  \item E.~Ghys,
        % L'invariant de Godbillon-Vey.
        \textit{S\'eminaire Bourbaki}, Vol.\ 1988/89,
        \textit{Ast\'erisque} \textbf{177--178} (1989), Exp.\ No.~706, 155--181.
  \item E.~Ghys,
        % Rigidite differentiable des groupes fuchsiens.
        \textit{Inst.\ Hautes \'Etudes Sci.\ Publ.\ Math.}\ \textbf{78} (1993), 163--185.
  \item E.~Ghys, R.~Langevin and P.~Walczak,
        % Entropie geometrique des feuilletages.
        \textit{Acta Math.}\ \textbf{160} (1988), no.~1--2, 105--142.
  \item 田村一郎，葉層のトポロジー，数学選書，岩波書店（1976）．
  \item P.~Tukia,
        % Homeomorphic conjugates of Fuchsian groups.
        \textit{J.\ Reine Angew.\ Math.}\ \textbf{391} (1988), 1--54.
  \item 矢野公一，力学系２，岩波講座現代数学の基礎，岩波書店（1998）．
\end{itemize}

\newpage

%======================================================================
% 講演２：松元重則
%======================================================================
\begin{center}
  {\large\bfseries 有界オイラー類と円周の同相群の共役}
\end{center}
\begin{flushright}
  \textbf{松元重則（日大理工）}
\end{flushright}

\medskip

種数 $g \geq 2$ の有向閉曲面 $\Sigma_g$ の基本群を $\Pi_g$ とし、
準同型 $\phi : \Pi_g \to \Diff$ を考える。
ここに $\Diff$ は $S^1$ の向きを保つ $C^r$ 級微分同相写像のなす群を表す。
古典的に $\phi$ には、Euler 数と呼ばれる数 $eu(\phi)$ が対応する。
これは $\phi$ から得られる $\Sigma_g$ 上の $S^1$ 束の Euler 類を
$\Sigma_g$ の基本類ではかった数であるが、また次のように与えることもできる。
曲面群 $\Pi_g$ の表示
\[
  \Pi_g = \langle a_1, b_1, \ldots, a_g, b_g \mid [a_1,b_1] \cdots [a_g,b_g] = 1 \rangle
\]
を（曲面の向きとあうように）とり、
準同型 $\phi$ による生成元の像の、被覆 $\R \to S^1$ に関する持ち上げを
$A_1, B_1, \ldots, A_g, B_g$ とする。
このとき元 $[A_1, B_1] \cdots [A_g, B_g]$ は恒等写像の持ち上げであり、
従って $\R$ 上の整数移動である。
この整数はすべてのものの選び方によらずに定まり、オイラー数 $eu(\phi)$ に一致する。

さて、あらゆる準同型 $\phi : \Pi_g \to \Diff$ について
不等式 $|eu(\phi)| \leq 2g - 2$ が成り立つことが知られている（Milnor--Wood の不等式）。
ちなみに曲面 $\Sigma_g$ を実現する Fuchs 群表示
$\Phi : \Pi_g \to PSL(2;\R) \subset \Diff$ をとれば $eu(\phi) = 2g - 2$ である。
ここに $PSL(2;\R)$ は無限円への作用により $\Diff$ の部分群とみなされる。
これら Fuchs 群表示はすべて位相共役であるが、
$C^r$ 共役（$r \geq 1$）においては Teichmüller 空間の点として分類される。

今回の講演で主にお話ししたいのは、次の定理である。

\medskip
\noindent
\textbf{定理１}\quad
準同型 $\phi : \Pi_g \to \Diff$（$r \geq 2$）で、等式 $eu(\phi) = 2g - 2$ を満たすものは
皆互いに位相共役である。

\medskip

この定理の証明には、M.~Gromov により、多様体のいわゆる単体的体積と関連して導入された
有界コホモロジーが、離散群の $S^1$ 上への作用を分類する手段として用いられる。

定理１は E.~Ghys により次の形に強められている。

\medskip
\noindent
\textbf{定理２}\quad
準同型 $\phi : \Pi_g \to \Diff$（$r \geq 3$）で、等式 $eu(\phi) = 2g - 2$ を満たすものは
Fuchs 群表示と $C^r$ 共役である。

\medskip

この証明は、力学系理論の様々な手法を応用しており大変おもしろい。
時間が許せばこれについても言及したい。

\newpage

%======================================================================
% 講演３：皆川宏之
%======================================================================
\begin{center}
  {\large\bfseries 双曲結晶群の $S^1$ の PL 同相群への表現：初等的な構成}
\end{center}
\begin{flushright}
  \textbf{皆川宏之（北大理）}
\end{flushright}

\medskip

$S^1$ の PL 同相写像群 $PL(S^1)$ を少しでも多くの人に知ってもらおうというのが今回の話の目的です。
この群への主なアプローチの仕方としては、
分類空間などを扱うホモトピー論的なもの、
有限表示部分群に対するデーン関数などを扱う群論的なもの、
そして $S^1$ への作用を扱う力学系的なものなどが挙げられます。
ここでは、力学系的アプローチを特に意識しながら
双曲結晶群すなわち Lie 群 $PSL(2,\mathbf{R})$ のココンパクト離散部分群から
$PL(S^1)$ への表現に関連するトピックスをいくつか紹介できればと考えています。

双曲結晶群は、双曲平面の理想境界への作用としてそれ自身 $S^1$ の実解析的微分同相群の部分群となりますが、
$PL(S^1)$ の部分群と位相共役となることがアノソフ流に対する手術を用いて示すことができます。
しかし、このようにして得られた PL 表現を実際に書き下したり解析したりすることは、
一般的には非常に大変な作業であり、さらにはこの手法では得られない PL 表現も数多く存在します。
そこで、講演では、
\begin{enumerate}
  \setlength{\itemsep}{2pt}
  \item[(1)] Ghys--Sergiescu による２進 PL 同相群のスムージング
  \item[(2)] アノソフ流の手術とアノソフ葉層の横断的 PL 構造
  \item[(3)] PL 表現の連続変形と離散的ゴドビヨン・ヴェイ不変量
\end{enumerate}
といった話題を念頭におきながら、曲面群・三角群といった双曲結晶群に対し
その PL 表現の初等的構成をお見せしたいと思っています。

\medskip

\begin{center}
  {\large 参考文献}
\end{center}

\begin{itemize}
  \setlength{\itemsep}{3pt}
  \item D.~Fried, \textit{Topology} \textbf{22} (1983), 299--303.
  \item E.~Ghys, \textit{Ann.\ Inst.\ Fourier} \textbf{37} (1987), 59--76.
  \item E.~Ghys, \textit{Topology} \textbf{26} (1987), 93--105.
  \item E.~Ghys and V.~Sergiescu,
        \textit{Comment.\ Math.\ Helv.}\ \textbf{62} (1987), 185--239.
  \item N.~Hashiguchi, \textit{Ann.\ Inst.\ Fourier} \textbf{42} (1992), 937--965.
  \item N.~Hashiguchi and H.~Minakawa,
        \textit{J.\ Fac.\ Sci.\ Univ.\ Tokyo Sect.\ IA Math.}\ \textbf{39} (1992), 271--278.
  \item H.~Minakawa, \textit{Topology} \textbf{30} (1991), 429--438.
  \item H.~Minakawa, \textit{Topology} \textbf{36} (1997), 775--781.
\end{itemize}

\end{document}